\documentclass[11pt,a4paper]{article}
% ============================================================
% Shared preamble for the MCA104 CNIA exam-preparation notes.
%
% libertinus and tcolorbox are installed in the user TeX tree
% (/home/elish4h/texmf) via:
%   tlmgr --usermode install --ignore-warning \
%     tcolorbox libertinus libertinus-type1 libertinus-otf \
%     libertinus-fonts libertinust1math
%
% The box styling matches the repository's notes-template.tex.
% ============================================================

\usepackage[margin=0.72in]{geometry}
\usepackage{libertinus}
\usepackage{microtype}
\usepackage[T1]{fontenc}
\usepackage{xcolor}
\usepackage{enumitem}
\usepackage{tabularx}
\usepackage{array}
\usepackage{booktabs}
\usepackage{longtable}
\usepackage{amsmath}
\usepackage{amssymb}
\usepackage{fancyhdr}
\usepackage{titlesec}
\usepackage[skins,breakable]{tcolorbox}
\usepackage{hyperref}
\usepackage{multicol}
\usepackage{parskip}

\definecolor{navy}{HTML}{183153}
\definecolor{blue}{HTML}{2563EB}
\definecolor{green}{HTML}{166534}
\definecolor{red}{HTML}{B91C1C}
\definecolor{gold}{HTML}{92400E}
\definecolor{lightblue}{HTML}{EFF6FF}
\definecolor{lightgreen}{HTML}{F0FDF4}
\definecolor{lightred}{HTML}{FEF2F2}
\definecolor{lightgold}{HTML}{FFFBEB}
\definecolor{grey}{HTML}{475569}

\hypersetup{colorlinks=true,linkcolor=blue,urlcolor=blue}

\titleformat{\section}{\Large\bfseries\color{navy}}{}{0pt}{}
\titleformat{\subsection}{\large\bfseries\color{navy}}{}{0pt}{}
\titleformat{\subsubsection}{\normalsize\bfseries\color{blue}}{}{0pt}{}

\setlist[itemize]{leftmargin=*,nosep}
\setlist[enumerate]{leftmargin=*,nosep}

% Allow tiny line-length overruns to be absorbed rather than overfull.
\setlength{\emergencystretch}{2em}

\tcbset{
  parbox=false,
  enhanced,
  breakable,
  colbacktitle=blue,
  coltitle=white,
  fonttitle=\bfseries,
  boxrule=0.5pt,
  arc=1.5pt,
  left=6pt,right=6pt,top=4pt,bottom=4pt,
  before skip=6pt,after skip=6pt,
}

\newtcolorbox{exambox}{colback=lightblue,colframe=blue,title=Exam Focus}
\newtcolorbox{rememberbox}{colback=lightgold,colframe=gold,title=Remember}
\newtcolorbox{dangerbox}{colback=lightred,colframe=red,title=Do Not Confuse}
\newtcolorbox{answerbox}{colback=lightgreen,colframe=green,title=Answer Skeleton}
\newtcolorbox{pyqbox}{colback=white,colframe=grey,title=Past-paper question}
\newtcolorbox{solbox}{colback=lightgreen,colframe=green,title=Solution}

\newcommand{\mkmarks}[1]{\hfill\textbf{[#1 marks]}}
\newcommand{\cmd}[1]{\texttt{#1}}
\newcommand{\F}[1]{\ensuremath{#1}}
\newcommand{\subject}[1]{\lhead{#1}\rhead{MCA104 Exam Notes}\cfoot{\thepage}}

\subject{MCA104 -- Module 1: Data Communication, Models and the Physical and Data Link Layers}

\title{\textbf{Module 1: Foundations, Models, Physical and Data Link Layers}\\}
\author{\large Data communication, OSI and TCP/IP, transmission media, switching, multiplexing,\\
line coding and PCM, error detection and correction, and the MAC sublayer}
\date{Exam: Wednesday 18 November 2026}

\begin{document}
\maketitle

\begin{exambox}
\textbf{Scope.} Syllabus 1.1 to 1.4: overview of data communication and networking; network models (OSI and TCP/IP); the physical layer (transmission media, switching, multiplexing); and the data link layer (error detection and correction, data link control).

Module 1 carries CO1 and dominates the CAE-1 paper. In the six papers we have, Module 1 accounts for \textbf{more marks than any other module}: every CAE-1 is almost entirely Module 1. Expect the layered-approach question, the OSI versus TCP/IP comparison, transmission media, multiplexing, switching, and a CRC calculation.
\end{exambox}

\tableofcontents
\newpage

\section{Data communication}

\subsection{Definition and components}

\textbf{Data communication} is the exchange of data between two devices through some transmission medium, in which the data is represented as agreed patterns of electromagnetic signals.

\begin{pyqbox}
\textbf{Sept 2026 CAE-1 Q1B(a).} Define Data Communication and Networking. Explain the five basic components of a data communication system.
\end{pyqbox}

\begin{solbox}
A data communication system has \textbf{five} components. Learn these five and the order, because the question is worth 10 marks and is pure recall.

\begin{center}\small
\begin{tabularx}{\textwidth}{>{\bfseries}p{0.18\textwidth} p{0.24\textwidth} X}
\toprule Component & Also called & Role\\
\midrule
Message & Data & The information to be communicated: text, numbers, images, audio, video.\\
Sender & Source & The device that sends the message: computer, phone, camera.\\
Receiver & Destination & The device that receives the message.\\
Transmission medium & Channel & The physical path the message travels: twisted pair, coaxial cable, fibre, radio waves.\\
Protocol & Rules & The set of rules governing communication: what is sent, how, and when. Without a protocol two devices may be connected but cannot communicate.\\
\bottomrule
\end{tabularx}
\end{center}

\textbf{Networking} is the interconnection of two or more devices so that they can share resources and exchange data. A \textbf{protocol} is the agreement that makes the exchange meaningful: syntax (format), semantics (meaning), and timing (when to send and how fast).
\end{solbox}

\subsection{Performance characteristics}

\begin{pyqbox}
\textbf{Sept 2026 CAE-1 Q1B(b).} Explain the important network performance characteristics: Bandwidth, Data Rate, Throughput, Latency, and Jitter.
\end{pyqbox}

\begin{solbox}
\begin{center}\small
\begin{tabularx}{\textwidth}{>{\bfseries}p{0.20\textwidth} X p{0.28\textwidth}}
\toprule Term & Meaning & Units and note\\
\midrule
Bandwidth & The \emph{capacity} of a medium to carry data: the range of frequencies it can pass. In digital terms, the maximum bit rate a link can support. & Hertz (analogue); bits per second (digital).\\
Data rate & The actual number of bits transmitted per second by a device. & bps, kbps, Mbps, Gbps.\\
Throughput & The \emph{actual} rate of successful data delivery, after protocol overhead, retransmissions and congestion. & Always $\le$ bandwidth.\\
Latency & Total time for a message to travel from sender to receiver: propagation $+$ transmission $+$ queueing $+$ processing. & Milliseconds.\\
Jitter & Variation in latency between successive packets. & Milliseconds. Fatal for real-time voice and video.\\
\bottomrule
\end{tabularx}
\end{center}

\textbf{Bandwidth versus throughput} is the distinction most often asked. Bandwidth is the pipe's size; throughput is what actually flows through it. A 100 Mbps link delivering 60 Mbps of useful data has a throughput of 60 Mbps.

\textbf{Latency versus jitter.} Latency is how late a packet is. Jitter is how \emph{unevenly late} successive packets are. A voice call tolerates a steady 150 ms delay; it does not tolerate 20 ms, 300 ms, 30 ms, 400 ms.
\end{solbox}

\begin{dangerbox}
\textbf{Bandwidth and data rate are sometimes used interchangeably, and sometimes not.} In strict usage, bandwidth is the frequency range of the medium, and data rate is the bit rate achieved. In the Shannon and Nyquist formulas the two are linked but distinct. If a question uses them loosely, define both terms in your answer and state the relation, and you will not lose marks either way.
\end{dangerbox}

\subsection{Transmission modes}

\begin{center}\small
\begin{tabularx}{\textwidth}{>{\bfseries}p{0.20\textwidth} X}
\toprule Mode & Description\\
\midrule
Simplex & One direction only, always. Keyboard to computer; radio broadcast.\\
Half duplex & Both directions, but only one at a time. Walkie-talkie.\\
Full duplex & Both directions simultaneously. Telephone; a normal switched Ethernet link.\\
\bottomrule
\end{tabularx}
\end{center}

\section{Network types and topologies}

\subsection{By geographical scope}

\begin{center}\small
\begin{tabularx}{\textwidth}{>{\bfseries}p{0.14\textwidth} p{0.20\textwidth} X}
\toprule Type & Scope & Characteristics\\
\midrule
PAN & A few metres & Personal area network: Bluetooth, USB, a phone tethered to a laptop.\\
LAN & A building or campus & High data rate, low cost, privately owned. Ethernet and Wi-Fi.\\
MAN & A city & Connects several LANs across a city; usually owned by a provider or a municipal body.\\
WAN & Country or global & Connects LANs and MANs over long distance; uses leased lines, MPLS, satellite. The Internet is the largest WAN.\\
\bottomrule
\end{tabularx}
\end{center}

\subsection{By topology}

\begin{center}\small
\begin{tabularx}{\textwidth}{>{\bfseries}p{0.13\textwidth} X p{0.30\textwidth}}
\toprule Topology & How it works & Strength / weakness\\
\midrule
Bus & All nodes tap one shared backbone cable. & Cheapest cabling; a break in the backbone kills the network; collisions.\\
Star & All nodes connect to a central hub or switch. & Easy to diagnose, one cable per node; the centre is a single point of failure.\\
Ring & Each node connects to two neighbours, forming a loop. & Predictable access (token ring); a break can isolate nodes.\\
Mesh & Every node connects to every other (full) or to several (partial). & Fault tolerant, many paths; expensive cabling, $n(n-1)/2$ links for full mesh.\\
Tree & A hierarchy of star-connected segments. & Scales to large sites; the root is critical.\\
Hybrid & A combination, as in most real networks. & Flexible; complex to manage.\\
\bottomrule
\end{tabularx}
\end{center}

\begin{rememberbox}
For a \textbf{full mesh} of $n$ nodes the number of links is $\dfrac{n(n-1)}{2}$ and each node needs $n-1$ interfaces. For 10 nodes that is 45 links. This arithmetic is a common 3-mark question.
\end{rememberbox}

\section{Network models}

\subsection{Why layer the network at all}

\begin{pyqbox}
\textbf{Sept 2026 CAE-1 Q1A.} Explain the importance of the layered approach in computer networking. Discuss how layering simplifies network design, implementation, troubleshooting, and interoperability, with suitable examples.
\end{pyqbox}

\begin{solbox}
Layering divides the complex job of moving data between applications on different machines into a stack of self-contained functions, where each layer uses the services of the layer below and provides services to the layer above.

\begin{enumerate}
\item \textbf{Divide and conquer.} Each layer solves one well-defined problem. The designer of the transport layer does not need to know whether the link is fibre or radio.
\item \textbf{Abstraction and independent implementation.} A layer's internals can be replaced without touching its neighbours. Ethernet can be swapped for Wi-Fi underneath IP and IP does not change.
\item \textbf{Interoperability through standard interfaces.} Because the boundary between layer $n$ and layer $n+1$ is standardised, a vendor can build a TCP stack and another vendor an Ethernet card, and they work together. This is what makes multi-vendor networking possible.
\item \textbf{Modular troubleshooting.} A failure can be localised to a layer. ``The cable is fine (layer 1), the switch sees the MAC (layer 2), but no ping reply (layer 3)'' is a systematic diagnosis. That is exactly what \texttt{ping}, \texttt{traceroute} and packet captures exploit.
\item \textbf{Reuse and specialisation.} The same transport layer serves web, mail and file transfer. Each new application does not re-implement reliability.
\item \textbf{Easier evolution.} IPv6 could be introduced at layer 3 while layers 1, 2 and 4 stayed as they were.
\end{enumerate}

\textbf{The cost:} layering adds header overhead and some duplication of function (error checking appears at several layers), and a strictly layered design can be slower than a monolithic one.
\end{solbox}

\subsection{The OSI reference model}

Seven layers, bottom to top. Learn the function \emph{and} the PDU \emph{and} the addressing used at each layer.

\begin{center}\small
\begin{tabularx}{\textwidth}{>{\bfseries}p{0.085\textwidth} p{0.13\textwidth} p{0.12\textwidth} X}
\toprule No. & Layer & PDU & Main function\\
\midrule
7 & Application & Data & Provides network services to the user: HTTP, FTP, SMTP, DNS.\\
6 & Presentation & Data & Translation, compression, encryption: character encoding, TLS.\\
5 & Session & Data & Establishes, manages and terminates sessions; synchronisation and checkpointing.\\
4 & Transport & Segment / Datagram & End-to-end delivery, segmentation and reassembly, flow and error control, port addressing. TCP and UDP.\\
3 & Network & Packet & Logical addressing and routing between networks. IP, ICMP, routing protocols.\\
2 & Data link & Frame & Framing, physical (MAC) addressing, error detection, flow control, media access. Ethernet, PPP.\\
1 & Physical & Bit & Transmission of raw bits over a medium: signalling, voltages, connectors, bit rate.\\
\bottomrule
\end{tabularx}
\end{center}

\begin{rememberbox}
\textbf{Mnemonic, bottom to top:} \emph{Please Do Not Throw Sausage Pizza Away} = Physical, Data link, Network, Transport, Session, Presentation, Application.

\textbf{Which layer, quickly:}
\begin{itemize}
\item ``Which route?'' \textrightarrow{} network (3)
\item ``Which port?'' \textrightarrow{} transport (4)
\item ``Which MAC address?'' \textrightarrow{} data link (2)
\item ``Which voltage, which pin?'' \textrightarrow{} physical (1)
\item ``Encrypt, compress, convert character set'' \textrightarrow{} presentation (6)
\end{itemize}
\end{rememberbox}

\subsection{The TCP/IP model}

\begin{pyqbox}
\textbf{Sept 2026 CAE-1 Q3B.} Explain the TCP/IP protocol suite. Discuss its layered architecture and compare it with the OSI reference model.
\end{pyqbox}

\begin{solbox}
The TCP/IP suite has \textbf{four} layers in the classic form, or five if the physical layer is counted separately.

\begin{center}\small
\begin{tabularx}{\textwidth}{>{\bfseries}p{0.20\textwidth} X p{0.24\textwidth}}
\toprule TCP/IP layer & Function & Protocols\\
\midrule
Application & Everything the user sees; combines OSI 5, 6 and 7 & HTTP, FTP, SMTP, DNS, DHCP, SNMP, TLS\\
Transport & End-to-end delivery, reliability, flow and congestion control & TCP, UDP, SCTP\\
Internet & Logical addressing and routing & IP, ICMP, IGMP, ARP (at the boundary), OSPF, BGP\\
Network access (link) & Framing, MAC addressing, media access, bit transmission & Ethernet, Wi-Fi, PPP\\
\bottomrule
\end{tabularx}
\end{center}

\textbf{Comparison with OSI:}

\begin{center}\small
\begin{tabularx}{\textwidth}{>{\bfseries}p{0.22\textwidth} X X}
\toprule Basis & OSI & TCP/IP\\
\midrule
Layers & 7 & 4\\
Nature & A theoretical reference model, model first then protocols & A practical model, protocols first then the description\\
Session and presentation & Separate layers & Folded into the application layer\\
Guaranteed delivery & Defined at the transport and network layers & Only the transport layer (TCP) guarantees it\\
Use & Teaching and diagnosis & Actual implementation on every host\\
Design origin & ISO, top-down & ARPANET/DoD, bottom-up\\
\bottomrule
\end{tabularx}
\end{center}

\textbf{Why TCP/IP won:} it was implemented and shipped before OSI was finished, it was free and open, and it was already running the Internet by the time OSI was standardised.
\end{solbox}

\subsection{Encapsulation and decapsulation}

\begin{pyqbox}
\textbf{Sept 2024 CAE-1 Q2.} Write notes on: 1.~Multiplexing 2.~Hierarchical Layering 3.~Encapsulation/Decapsulation 4.~PDU.
\end{pyqbox}

\begin{solbox}
\textbf{Encapsulation} is the process of adding a header (and sometimes a trailer) at each layer as data moves \emph{down} the stack.

\begin{center}\small
\begin{tabularx}{\textwidth}{>{\bfseries}p{0.22\textwidth} X}
\toprule Layer & What is added\\
\midrule
Application & The user's message (e.g.\ an HTTP request)\\
Transport & TCP header with source and destination port, sequence number, checksum\\
Network & IP header with source and destination IP, TTL, protocol\\
Data link & Frame header with source and destination MAC, plus a trailer with the FCS\\
Physical & Converts the frame into signals (bits)\\
\bottomrule
\end{tabularx}
\end{center}

\textbf{Decapsulation} is the reverse: each layer strips and inspects its own header and passes the payload up.

\textbf{PDU (protocol data unit)} is the name of the data unit at each layer: \textbf{bit} at physical, \textbf{frame} at data link, \textbf{packet} at network, \textbf{segment} at transport (datagram for UDP), and \textbf{message} or \textbf{data} above that.

\textbf{Hierarchical layering} is the arrangement itself: a stack of layers in which each layer offers services to the one above and consumes services from the one below, communicating logically with its peer layer in the other host through the headers it adds.
\end{solbox}

\section{Physical layer: transmission media}

\begin{pyqbox}
\textbf{Sept 2026 CAE-1 Q2B.} Describe the different transmission media used in computer networks. Explain guided and unguided media with their advantages, disadvantages, and applications.
\end{pyqbox}

\begin{solbox}
Transmission media divide into \textbf{guided} (the signal is confined to a physical path) and \textbf{unguided} (the signal radiates through free space).

\textbf{Guided media:}

\begin{center}\small
\begin{tabularx}{\textwidth}{>{\bfseries}p{0.17\textwidth} X p{0.24\textwidth}}
\toprule Medium & Characteristics & Use\\
\midrule
Twisted pair (UTP/STP) & Two insulated copper wires twisted to cancel interference. UTP categories 5e, 6, 6a, 7. Cheap and easy to install; limited distance ($\approx$100 m) and bandwidth. & LAN cabling, telephone lines.\\
Coaxial cable & A central conductor surrounded by a shield. Higher bandwidth and better shielding than twisted pair. Bulky and expensive. & Cable TV, older Ethernet (10Base2), broadband.\\
Optical fibre & Glass or plastic core carrying light. Single-mode for long haul, multimode for short. Enormous bandwidth, immune to electromagnetic interference, very secure, low attenuation; expensive and fragile. & Backbones, undersea cables, FTTH, data centres.\\
\bottomrule
\end{tabularx}
\end{center}

\textbf{Unguided media:}

\begin{center}\small
\begin{tabularx}{\textwidth}{>{\bfseries}p{0.17\textwidth} X p{0.24\textwidth}}
\toprule Medium & Characteristics & Use\\
\midrule
Radio waves & Omni-directional, penetrates walls, 3 kHz to 1 GHz. & AM/FM radio, Wi-Fi, mobile.\\
Microwaves & Line of sight, high frequency, needs dish alignment and towers; affected by rain. & Point-to-point links, satellite.\\
Infrared & Very short range, cannot pass through walls, needs line of sight. & TV remotes, short-range device links.\\
\bottomrule
\end{tabularx}
\end{center}

\textbf{Choosing a medium} comes down to four trade-offs: bandwidth needed, distance, immunity to interference, and cost. Fibre wins on all but cost and fragility, which is why it is used for the backbone and copper for the last hundred metres.
\end{solbox}

\section{Switching}

\begin{pyqbox}
\textbf{Sept 2026 CAE-1 Q3A.} Explain different switching techniques used in computer networks. Compare circuit switching and packet switching.
\end{pyqbox}

\begin{solbox}
\textbf{Switching} is the technique by which a switched network moves data from source to destination through intermediate nodes.

\begin{center}\small
\begin{tabularx}{\textwidth}{>{\bfseries}p{0.18\textwidth} X}
\toprule Technique & How it works\\
\midrule
Circuit switching & A dedicated physical path is established end to end before any data flows, and it is held for the whole session. Resources are reserved.\\
Message switching & The whole message is stored at each intermediate node, then forwarded. No dedicated path; store-and-forward. Obsolete.\\
Packet switching & The message is broken into packets that are forwarded independently, each carrying the destination address, sharing the links statistically.\\
\bottomrule
\end{tabularx}
\end{center}

\begin{center}\small
\begin{tabularx}{\textwidth}{>{\bfseries}p{0.22\textwidth} X X}
\toprule Basis & Circuit switching & Packet switching\\
\midrule
Path & Dedicated, established first & No dedicated path; each packet routed independently\\
Resource use & Reserved even when idle (wasteful) & Shared on demand (efficient)\\
Set-up delay & Yes, before data & None\\
Delay & Constant after set-up & Variable, due to queueing\\
Congestion & Refused at set-up if no capacity & Packets queue and may be dropped\\
Overhead & Minimal & Per-packet header\\
Best for & Voice, real-time streams & Burst data: web, email, file transfer\\
Example & PSTN & The Internet\\
\bottomrule
\end{tabularx}
\end{center}

\textbf{Why the Internet uses packet switching:} data traffic is bursty. Reserving a circuit for a user who is reading a web page wastes capacity that another user could use. Statistical multiplexing gives far higher utilisation at the cost of variable delay.
\end{solbox}

\section{Multiplexing}

\begin{pyqbox}
\textbf{Sept 2026 CAE-1 Q4A and Oct 2025 Q1(a).} What is multiplexing? Explain the need for multiplexing. Discuss FDM, TDM and WDM with suitable examples.
\end{pyqbox}

\begin{solbox}
\textbf{Multiplexing} is the technique of combining several signals onto one high-capacity link so that it can be shared, and separating them again at the receiver (\textbf{demultiplexing}). The device is a \textbf{MUX}; the reverse is a \textbf{DEMUX}.

\textbf{The need:} a single link almost always has far more capacity than one user's traffic, and laying a separate cable per user is prohibitively expensive. Multiplexing converts one high-bandwidth link into many logical channels, so cost per user falls and utilisation rises.

\begin{center}\small
\begin{tabularx}{\textwidth}{>{\bfseries}p{0.16\textwidth} p{0.19\textwidth} X}
\toprule Technique & How it separates users & Notes\\
\midrule
FDM & Different \textbf{frequency} bands & Analogue; each user gets a slice of the spectrum continuously. Used in radio and TV broadcast, and ADSL.\\
TDM & Different \textbf{time slots} & Digital; users take turns on the whole link. Synchronous TDM assigns fixed slots (wasteful when a user is idle); statistical TDM assigns slots on demand.\\
WDM & Different \textbf{wavelengths} (colours) of light & The optical version of FDM, used in fibre backbones. DWDM packs dozens of wavelengths into one fibre.\\
CDM/CDMA & Different \textbf{codes} & All users transmit at once on the whole band; each uses a unique spreading code. Used in 3G mobile.\\
\bottomrule
\end{tabularx}
\end{center}

\begin{rememberbox}
\textbf{FDM vs TDM in one line each.} FDM: every user transmits all the time, on a narrow band. TDM: every user uses the whole band, for a short slice of time. If the question asks for a difference, this sentence plus a diagram is a full answer.
\end{rememberbox}

\textbf{The hierarchy:} TDM carriers are built in steps. In the North American hierarchy, 24 DS-0 channels (64 kbps each) form a DS-1/T-1 at 1.544 Mbps; in the European hierarchy, 30 channels plus 2 for signalling form an E-1 at 2.048 Mbps.
\end{solbox}

\section{Digital transmission: line coding and PCM}

\begin{pyqbox}
\textbf{Sept 2024 CAE-1 Q1.} Explain the polar coding schemes of digital signals.
\end{pyqbox}

\begin{solbox}
\textbf{Line coding} converts a sequence of bits into a digital signal. In \textbf{polar} schemes, two voltage levels are used, one positive and one negative, so the average DC level is zero.

\begin{center}\small
\begin{tabularx}{\textwidth}{>{\bfseries}p{0.20\textwidth} p{0.22\textwidth} X}
\toprule Scheme & Rule & Property\\
\midrule
NRZ-L & Level depends on the bit: one level for 0, the other for 1 & Simple; long runs of the same bit cause loss of synchronisation\\
NRZ-I & Invert the level on a 1, hold it on a 0 & Solves the long-run problem for 1s, not for 0s\\
Manchester & Transition in the middle of every bit; low-to-high for 0, high-to-low for 1 & Self-clocking: the clock is recovered from the signal. Used in classic Ethernet\\
Differential Manchester & Always a transition in the middle; a transition at the start of the bit means 0, none means 1 & Self-clocking, and polarity-independent\\
AMI & 0 is zero volts; 1s alternate between + and -- & Good DC balance; used in T1 lines\\
\bottomrule
\end{tabularx}
\end{center}

The cost of Manchester coding is that it needs twice the bandwidth of NRZ for the same bit rate, because every bit carries a transition. This is the standard trade-off: self-clocking costs bandwidth.
\end{solbox}

\begin{pyqbox}
\textbf{Sept 2024 CAE-1 Q3.} What is Pulse Code Modulation (PCM), and how does it convert analog signals into digital form for transmission?
\end{pyqbox}

\begin{solbox}
\textbf{PCM} is the standard method of digitising an analogue signal. Three steps:

\begin{enumerate}
\item \textbf{Sampling.} Measure the analogue signal at regular intervals. By the \textbf{Nyquist theorem}, the sampling rate must be at least twice the highest frequency present: $f_s \ge 2 f_{max}$. For voice band-limited to 4 kHz, $f_s = 8000$ samples per second.
\item \textbf{Quantisation.} Round each sample to the nearest of $L = 2^n$ discrete levels. The rounding error is \textbf{quantisation noise}; using more levels reduces it. Non-uniform quantisation (companding, as in $\mu$-law) gives small signals finer resolution.
\item \textbf{Encoding.} Assign a binary code of $n$ bits to each level.
\end{enumerate}

\textbf{Bit rate:} $R = f_s \times n$ bits per second. For digitised voice, $f_s = 8000$ and $n = 8$, so $R = 64$ kbps. That 64 kbps is the DS-0 channel on which all TDM hierarchies are built.

\textbf{Quantisation error} is at most half a step: SNR$_{dB} \approx 6.02n + 1.76$ for an $n$-bit quantiser. Each extra bit buys about 6 dB.
\end{solbox}

\begin{rememberbox}
\textbf{Two capacity formulas.} Nyquist for a noiseless channel: $C = 2B\log_2 L$ bps, where $B$ is bandwidth and $L$ the number of signal levels. Shannon for a noisy channel: $C = B\log_2(1 + S/N)$. Shannon is the hard ceiling; Nyquist is the ceiling when there is no noise.
\end{rememberbox}

\section{Data link layer}

\subsection{Functions}

\begin{pyqbox}
\textbf{Oct 2025 Q1(a)(iii).} Describe the concept of framing in data link control.
\end{pyqbox}

\begin{solbox}
The data link layer turns the physical layer's raw bit stream into manageable units and manages the link. Its functions are:

\begin{itemize}
\item \textbf{Framing}: dividing the bit stream into frames, each with a header, payload and trailer.
\item \textbf{Physical addressing}: source and destination MAC addresses.
\item \textbf{Error control}: detect (and sometimes correct) errors with a trailer checksum.
\item \textbf{Flow control}: prevent a fast sender from overwhelming a slow receiver.
\item \textbf{Media access control}: decide who may transmit on a shared medium.
\end{itemize}

\textbf{Framing} is necessary because the physical layer delivers an undifferentiated stream of bits. The receiver must know where a frame begins and ends. Three techniques:

\begin{center}\small
\begin{tabularx}{\textwidth}{>{\bfseries}p{0.20\textwidth} X}
\toprule Technique & How\\
\midrule
Character count & A field in the header states the frame length. Fragile: one corrupted count desynchronises everything.\\
Byte stuffing & A special flag byte marks the start and end; any flag byte occurring in the data is escaped with an extra byte.\\
Bit stuffing & A flag pattern (01111110) delimits frames; in the data, a 0 is inserted after any five consecutive 1s. HDLC uses this.\\
\bottomrule
\end{tabularx}
\end{center}
\end{solbox}

\subsection{Error detection: parity, checksum and CRC}

\subsubsection{Simple and two-dimensional parity}

A \textbf{single parity bit} makes the number of 1s even (even parity) or odd. It detects any \emph{odd} number of bit errors, so it misses all two-bit errors.

\textbf{Two-dimensional parity} arranges data as a block and computes parity for each row \emph{and} each column. It detects all single, double and triple errors, and can \emph{correct} a single error by locating the failing row and column intersection.

\subsubsection{Internet checksum}

Send the 16-bit one's-complement sum of the data words as a checksum; the receiver adds all words including the checksum and expects all ones.

\begin{exambox}
\textbf{Worked.} Words \texttt{1001 1001 1001 1001} and \texttt{0101 0101 0101 0101}.

Sum $= 1110\,1110\,1110\,1110$; checksum $=$ one's complement $= \texttt{0001 0001 0001 0001}$.

Receiver adds all three: $1111\,1111\,1111\,1111$ --- all ones, so no error is detected. A checksum is weaker than CRC; it is used in IP and UDP because it is cheap in software.
\end{exambox}

\subsubsection{CRC --- the calculation you will be asked to do}

\begin{pyqbox}
\textbf{Sept 2026 CAE-1 Q4B.} Given the data bits 1101011011 and the generator bits 10011: (a) append the required number of zeros; (b) perform the modulo-2 binary division, showing all steps; (c) determine the CRC remainder and construct the final transmitted codeword.
\end{pyqbox}

\begin{solbox}
\textbf{The rule:} if the generator has $n+1$ bits, append $n$ zeros, divide modulo 2, and replace the zeros with the remainder.

Generator \texttt{10011} has 5 bits, so $n = 4$ and we append \textbf{4 zeros}.

\textbf{Step 1 --- append:}
\[
\texttt{1101011011} + \texttt{0000} = \texttt{11010110110000}
\]

\textbf{Step 2 --- divide modulo 2} (XOR, no carries). Each time the leading bit is 1, XOR the generator into the window and drop the leading 0; when it is 0, shift one place.

\begin{center}\small
\begin{tabular}{r|l}
\toprule Step & Current window\\
\midrule
start & \texttt{11010110110000}\\
XOR at position 0 & \texttt{10011} $\Rightarrow$ \texttt{01001110110000}\\
shift to position 1 & \texttt{1001110110000}\\
XOR at position 1 & \texttt{10011} $\Rightarrow$ \texttt{0000010110000}\\
shift to position 6 & \texttt{10110000}\\
XOR at position 6 & \texttt{10011} $\Rightarrow$ \texttt{0010100}\\
shift to position 8 & \texttt{10100}\\
XOR at position 8 & \texttt{10011} $\Rightarrow$ \texttt{00111}\\
remainder & \texttt{1110}\\
\bottomrule
\end{tabular}
\end{center}

Let us verify the arithmetic rather than trust it. Position 0: \texttt{11010} $\oplus$ \texttt{10011} = \texttt{01001}. Position 1: window \texttt{10011} $\oplus$ \texttt{10011} = \texttt{00000}. What remains after the two leading XORs is the low bits \texttt{10110000}. Position 6: \texttt{10110} $\oplus$ \texttt{10011} = \texttt{00101}, leaving \texttt{0100}. There is no 1 left at a position that can start a division except the last window \texttt{10100} at position 8: \texttt{10100} $\oplus$ \texttt{10011} = \texttt{00111}. The remainder is \texttt{1110}.

\textbf{Step 3 --- the codeword:} replace the four appended zeros with the remainder \texttt{1110}:
\[
\textbf{11010110111110}
\]

\textbf{Step 4 --- why it works.} The receiver divides the received 14 bits by the same generator. If the remainder is \texttt{0000} the frame is accepted; anything else is discarded. This works because the transmitted polynomial is exactly divisible by $G(x)$.
\end{solbox}

\begin{pyqbox}
\textbf{Oct 2025 CAE-II Q1A and Oct 2024 CAE-II Q1A (the same question in both papers).} A bit stream 10011101 is transmitted using the standard CRC method. The generator polynomial is $x^3+1$. (a) What is the actual bit string transmitted? (b) Suppose the third bit from the left is inverted during transmission. How will the receiver detect this error?
\end{pyqbox}

\begin{solbox}
\textbf{Step 1 --- the generator.} $G(x) = x^3 + 1$ means coefficients at $x^3$ and $x^0$, so the bit pattern is \texttt{1001} (4 bits, degree 3). Append \textbf{3} zeros.

\[
\texttt{10011101} + \texttt{000} = \texttt{10011101000}
\]

\textbf{Step 2 --- divide:}

\begin{center}\small
\begin{tabular}{r|l}
\toprule Step & Window\\
\midrule
start & \texttt{10011101000}\\
XOR at position 0 & \texttt{1001} $\Rightarrow$ \texttt{00001101000}\\
shift to position 4 & \texttt{1101000}\\
XOR at position 4 & \texttt{1001} $\Rightarrow$ \texttt{0100000}\\
shift to position 5 & \texttt{100000}\\
XOR at position 5 & \texttt{1001} $\Rightarrow$ \texttt{000100}\\
remainder & \texttt{100}\\
\bottomrule
\end{tabular}
\end{center}

The remainder is \texttt{100}, so the transmitted string is
\[
\textbf{10011101100}
\]

\textbf{Step 3 --- the inverted bit.} The third bit from the left of \texttt{10011101100} is the 0 in position 3 (positions 1, 2, 3 are 1, 0, 0), so it becomes 1:
\[
\text{received} = \texttt{10111101100}
\]

Divide the received bits by \texttt{1001}. The remainder is \texttt{100}, which is \textbf{not zero}, so the receiver discards the frame and requests a retransmission.

\textbf{The point to write down:} the receiver does not need to know \emph{where} the error is or \emph{what} it was. A non-zero remainder is enough. That is why CRC is the standard link-layer detector.
\end{solbox}

\begin{dangerbox}
\textbf{Modulo-2 arithmetic means XOR, not subtraction.} There are no borrows and no carries. If you find yourself writing a 2 or a --1 in the division, stop: you are doing ordinary binary subtraction.
\end{dangerbox}

\subsection{Error correction: Hamming code}

To correct a single error you need enough redundancy to identify \emph{which} bit. With $m$ data bits and $r$ check bits, the condition is
\[
2^r \ge m + r + 1.
\]
For $m = 4$, $r = 3$ satisfies $8 \ge 8$, giving the classic (7,4) Hamming code. Check bits sit at positions 1, 2, 4, 8, \dots; each check bit covers the positions whose binary index contains that bit.

\begin{rememberbox}
\textbf{Detection versus correction.} Detection only needs to answer ``did something go wrong''. Correction must answer ``what exactly went wrong''. Correction therefore needs far more redundancy, which is why links normally detect and retransmit (ARQ) rather than correct, and only real-time links (satellite, deep space) pay for forward error correction.
\end{rememberbox}

\subsection{Data link control: flow control and ARQ}

\begin{pyqbox}
\textbf{Oct 2025 Q1(b)(i).} Explain the working principle of the Go-Back-N ARQ protocol. Compare it with the Selective Repeat ARQ protocol in terms of efficiency.
\end{pyqbox}

\begin{solbox}
\textbf{Go-Back-N.} The sender may have up to $N$ unacknowledged frames outstanding. The receiver accepts only in-order frames and sends a cumulative ACK. If a frame is lost, the receiver discards everything after it, and the sender, on timeout, \textbf{retransmits that frame and all subsequent ones}.

Window size: $N = 2^m - 1$ for an $m$-bit sequence number.

\textbf{Selective Repeat.} Same windowing, but the receiver buffers out-of-order frames and ACKs each one individually; the sender retransmits \textbf{only} the missing frame.

Window size: $N = 2^{m-1}$.

\begin{center}\small
\begin{tabularx}{\textwidth}{>{\bfseries}p{0.22\textwidth} X X}
\toprule Basis & Go-Back-N & Selective Repeat\\
\midrule
Receiver buffer & None; out-of-order frames discarded & Buffers out-of-order frames\\
Retransmission & The lost frame and every frame after it & Only the lost frame\\
Window size & $2^m - 1$ & $2^{m-1}$\\
Efficiency on a lossy link & Falls sharply as the error rate rises & Stays high\\
Complexity & Simple sender and receiver & More complex receiver and buffering\\
\bottomrule
\end{tabularx}
\end{center}

\textbf{Efficiency.} For a link with transmission time $T_t$ and one-way propagation $T_p$, let $a = T_p / T_t$. The sender can keep the link busy only if
\[
N \ge 1 + 2a.
\]
The utilisation of Go-Back-N is $U = \dfrac{N}{1+2a}$ (capped at 1). For example with $a = 2$, a window of 5 gives $U = 5/5 = 1$, i.e.\ 100\%; a window of 3 gives $U = 3/5 = 0.6$, i.e.\ 60\%.

Selective Repeat reaches the same utilisation but wastes no capacity on retransmitting frames that arrived correctly. Go-Back-N is cheaper to implement; Selective Repeat is more efficient. That is the trade-off the question wants.
\end{solbox}

\begin{pyqbox}
\textbf{Oct 2024 CAE-II Q1B.} Explain the sliding window protocols in detail.
\end{pyqbox}

\begin{solbox}
A \textbf{sliding window} bounds how many frames may be in flight. The sender keeps a window of at most $N$ unacknowledged frames; each ACK slides the window forward. This gives flow control (the window never exceeds what the receiver can buffer) and efficiency (the sender does not stop and wait for each frame).

\begin{itemize}
\item Window size 1 with stop-and-wait: one frame at a time. Utilisation $= \dfrac{1}{1+2a}$, which collapses on links with a large propagation delay, such as satellite.
\item Window size $N$: pipelining. Utilisation $= \dfrac{N}{1+2a}$, capped at 1.
\end{itemize}

The two ARQ implementations of the sliding window are Go-Back-N and Selective Repeat, compared above. \textbf{Piggybacking} is the optimisation in which a data frame travelling in the reverse direction carries the ACK for the forward direction, saving a separate ACK frame.
\end{solbox}

\section{The MAC sublayer}

\begin{pyqbox}
\textbf{Sept 2026 CAE-1 Q2A.} Explain the concept of Random Access Protocols used in the MAC sublayer. Discuss the working principles of ALOHA, Slotted ALOHA, CSMA, and CSMA/CD, highlighting their differences and limitations.
\end{pyqbox}

\begin{solbox}
\textbf{Random access} means no station is assigned a turn in advance; any station may transmit when it has data, and a \textbf{contention resolution} procedure handles the resulting collisions. This is efficient for bursty traffic, which is why Ethernet uses it.

\begin{center}\small
\begin{tabularx}{\textwidth}{>{\bfseries}p{0.20\textwidth} X p{0.20\textwidth}}
\toprule Protocol & Principle & Limitation\\
\midrule
ALOHA & Transmit whenever you have data. If no ACK, wait a random time and retransmit. & Maximum throughput $18.4\%$ of the channel; collisions are frequent\\
Slotted ALOHA & Time is divided into slots; transmit only at a slot boundary & Doubles efficiency to $36.8\%$; still collision-prone\\
CSMA & \emph{Carrier sense}: listen before transmitting. Transmit only if the medium is idle & Collisions still possible because of propagation delay\\
CSMA/CD & CSMA plus \emph{collision detection}: keep listening while transmitting, abort on collision, send a jam signal, back off exponentially & Only works on wired media where a station can hear while it talks\\
CSMA/CA & Collision \emph{avoidance}: wait a random backoff before transmitting, and optionally use RTS/CTS to reserve the medium & Needed on wireless, where a station cannot detect collisions\\
\bottomrule
\end{tabularx}
\end{center}

\textbf{Why throughput is $18.4\%$ for ALOHA:} in pure ALOHA a frame is vulnerable for two frame times, so the probability of no collision is $e^{-2G}$ and the throughput $S = Ge^{-2G}$ peaks at $G = 0.5$, giving $S = 0.184$. For slotted ALOHA the vulnerable period halves, so $S = Ge^{-G}$ peaks at $G = 1$ with $S = 0.368$.
\end{solbox}

\begin{pyqbox}
\textbf{Oct 2025 CAE-II Q1B.} A LAN has 5 computers (A, B, C, D, E) connected using CSMA/CD. The time it takes a signal to travel between any two computers is 25 $\mu$s. Each frame takes 100 $\mu$s to send. Computer A starts sending a frame at the same time as computer C. (a) What happens when A and C send at the same time? (b) How do the computers know that a collision has happened? (c) What do the computers do after detecting the collision?
\end{pyqbox}

\begin{solbox}
\textbf{(a) What happens.} Both stations sense the medium idle and begin transmitting simultaneously. Their signals meet at the midpoint of the path between them, about 12.5 $\mu$s after both started, and the overlapping signals corrupt each other. Neither frame is delivered.

\textbf{(b) How they know.} A station transmitting with CSMA/CD continues to monitor the medium. When it senses a signal level or waveform different from its own, it concludes a collision has occurred. The time to notice is at most one round trip, $2T_p = 2 \times 25 = 50\ \mu$s.

\textbf{The key condition.} CSMA/CD only works if the station is \emph{still transmitting} when it learns of the collision, that is
\[
T_t \ge 2T_p.
\]
Here $T_t = 100\ \mu$s and $2T_p = 50\ \mu$s, so the condition holds comfortably: the station is only halfway through its frame when it detects the collision. If the frame were shorter than 50 $\mu$s, the station would finish and believe it had succeeded.

\textbf{(c) What they do next.} Each station aborts its transmission immediately and sends a short \textbf{jam signal} (32 to 48 bits) so that every other station also notices the collision. Then both enter \textbf{binary exponential backoff}: after the $n$-th collision a station waits a random number of \emph{slot times} ($2T_p = 50\ \mu$s) chosen from $\{0, 1, \dots, 2^n - 1\}$, and tries again. After 16 failed attempts the frame is abandoned and reported as an error.

\textbf{Efficiency of this LAN:} $U = \dfrac{T_t}{T_t + 2T_p} = \dfrac{100}{100+50} = 0.667$, that is 66.7\%.
\end{solbox}

\begin{rememberbox}
\textbf{CSMA/CD versus CSMA/CA.} Wired Ethernet can detect a collision because the transmitter can listen to the same cable it is transmitting on. A wireless station cannot: its own transmission drowns out everything else (the ``hidden node'' problem), so it cannot detect a collision and must \emph{avoid} one instead, using random backoff and optional RTS/CTS. That single physical fact is the whole reason 802.3 uses CD and 802.11 uses CA.
\end{rememberbox}

\section{Answer skeletons}

\begin{answerbox}
\textbf{Define-and-list (3 to 10 marks).} Give a one-sentence definition, then a labelled list. For the five components of data communication, name all five in order and add a clause each. Never leave a list unlabelled.
\end{answerbox}

\begin{answerbox}
\textbf{Compare-two-things (10 marks).} A table, at least five rows, each row a single basis of comparison. OSI versus TCP/IP; circuit versus packet; FDM versus TDM; GBN versus SR; CSMA/CD versus CSMA/CA; guided versus unguided media. Write the table first, then one sentence of conclusion.
\end{answerbox}

\begin{answerbox}
\textbf{CRC calculation (10 marks).} State the generator in bits, state how many zeros to append and why, show the division step by step, box the remainder, then write the codeword. The steps carry the marks, not the answer.
\end{answerbox}

\begin{answerbox}
\textbf{Explain-a-mechanism (10 marks).} Definition, then the mechanism in numbered steps, then a diagram, then one advantage and one limitation. For a protocol this is: what it does, how it does it, when it fails.
\end{answerbox}

\section{Self-test}

\begin{enumerate}
\item Define data communication and list its five components.
\item Differentiate between bandwidth, data rate and throughput.
\item What is jitter and why does it matter for voice?
\item State the seven OSI layers with their PDUs.
\item Compare OSI and TCP/IP in a table.
\item What is encapsulation? Name the PDU at each layer.
\item Distinguish guided from unguided media with two examples each.
\item Compare circuit switching and packet switching.
\item What is the need for multiplexing? Differentiate FDM and TDM.
\item What is PCM? State the three steps and the bit rate for digitised voice.
\item State the Nyquist and Shannon capacity formulas.
\item What is framing? Name three framing techniques.
\item Perform the CRC for data 10011101 with generator $x^3+1$.
\item How many check bits are needed to correct a single error in 4 data bits?
\item State the Go-Back-N and Selective Repeat window sizes for a 3-bit sequence number.
\item Derive the condition for CSMA/CD to work.
\item Why is the maximum throughput of ALOHA 18.4\%?
\item Why does wireless use CSMA/CA instead of CSMA/CD?
\end{enumerate}

\subsection*{Answers}
\begin{enumerate}
\item The exchange of data between devices over a medium. Components: message, sender, receiver, transmission medium, protocol.
\item Bandwidth is the capacity of the medium (Hz or max bit rate); data rate is the actual bit rate of a device; throughput is the successful delivery rate after overhead, always $\le$ bandwidth.
\item Jitter is the variation in latency between successive packets. Real-time audio and video need steady delay, so jitter causes glitches even when average latency is acceptable.
\item Physical (bit), data link (frame), network (packet), transport (segment), session, presentation, application (data).
\item See the comparison table: 7 layers versus 4; theoretical versus practical; session and presentation folded in; TCP/IP is what is deployed.
\item Adding a header at each layer on the way down. PDUs: bit, frame, packet, segment, data.
\item Guided confines the signal: twisted pair, coaxial, fibre. Unguided radiates: radio, microwave, infrared.
\item Circuit switching reserves a dedicated path and is good for voice; packet switching forwards independent packets and is good for bursty data.
\item Multiplexing shares one link among many users to cut cost and raise utilisation. FDM separates by frequency band; TDM separates by time slot.
\item PCM digitises analogue signals by sampling, quantising and encoding. For voice, $f_s = 8000$ and $n = 8$, giving 64 kbps.
\item Nyquist $C = 2B\log_2 L$; Shannon $C = B\log_2(1+S/N)$.
\item Framing divides the bit stream into frames. Techniques: character count, byte stuffing, bit stuffing.
\item Generator \texttt{1001}, append 3 zeros, remainder \texttt{100}, codeword \texttt{10011101100}.
\item $r = 3$, since $2^3 = 8 \ge 4 + 3 + 1 = 8$.
\item GBN: $2^3 - 1 = 7$. Selective Repeat: $2^{3-1} = 4$.
\item $T_t \ge 2T_p$: the frame must take at least as long as the round trip, otherwise the sender finishes before learning of the collision.
\item A frame is vulnerable for two frame times, so $S = Ge^{-2G}$, which peaks at $G = 0.5$ with $S = 0.184$.
\item A wireless station cannot hear collisions while transmitting because its own signal dominates, so it must avoid them by backoff and RTS/CTS.
\end{enumerate}

\section{Last-minute page}

\begin{multicols}{2}
\textbf{DATA COMM}\\
5 components: message, sender,\\
receiver, medium, protocol\\[3pt]
Bandwidth $>$ throughput always\\
Latency vs jitter\\[5pt]
\textbf{OSI 7}\\
Physical, Data link, Network,\\
Transport, Session, Presentation,\\
Application\\
PDU: bit, frame, packet, segment\\[5pt]
\textbf{TCP/IP 4}\\
Application, Transport,\\
Internet, Network access\\[5pt]
\textbf{MEDIA}\\
Guided: twisted pair, coax, fibre\\
Unguided: radio, microwave, IR\\[5pt]
\textbf{SWITCHING}\\
Circuit: reserved path, voice\\
Packet: independent, data\\
Message: store and forward\\[5pt]
\textbf{MULTIPLEXING}\\
FDM: frequency, all the time\\
TDM: time slot, whole band\\
WDM: wavelengths, fibre\\[5pt]
\textbf{PCM}\\
Sample $\to$ quantise $\to$ encode\\
$f_s \ge 2f_{max}$\\
8000 $\times$ 8 = 64 kbps\\[5pt]
\textbf{CAPACITY}\\
Nyquist $2B\log_2 L$\\
Shannon $B\log_2(1+S/N)$\\[5pt]
\textbf{ERROR CONTROL}\\
Parity: odd/even, misses 2-bit\\
Checksum: 1s complement sum\\
CRC: XOR division, remainder\\[5pt]
\textbf{ARQ}\\
GBN window $2^m-1$\\
SR window $2^{m-1}$\\
$N \ge 1+2a$ for full use\\[5pt]
\textbf{MAC}\\
ALOHA 18.4\%, slotted 36.8\%\\
CSMA/CD: wired, $T_t \ge 2T_p$\\
CSMA/CA: wireless, backoff\\[5pt]
\textbf{DO NOT CONFUSE}\\
Bandwidth vs throughput\\
Latency vs jitter\\
FDM (frequency) vs TDM (time)\\
CSMA/CD (detect) vs CSMA/CA (avoid)
\end{multicols}

\vfill
\begin{center}\textit{Module 1 is definition-and-table territory. Learn the five components, the seven layers with PDUs, and the three comparisons: OSI/TCP-IP, circuit/packet, FDM/TDM.}\end{center}

\end{document}
